\section{Desarrollo del modelo Canvas}

\subsection{Segmentos de mercado}

Pulso se dirige principalmente a dueños y administradores que necesitan ordenar
los cobros, las membresías y la asistencia. También puede servir a entrenadores
personales que organizan planes y dan seguimiento a varios clientes. Los
miembros usarán la plataforma para consultar su información; su experiencia
será importante para que el gimnasio decida seguir con el servicio.

La primera prueba se propone en Pérez Zeledón, en un gimnasio administrado por
una o pocas personas, con acceso a internet y disposición para recibir apoyo
durante el piloto. La cercanía facilita las visitas. Más adelante se podrían
atender gimnasios de otros lugares mediante el acceso en línea.

\begin{table}[htbp]
\caption{Segmentos y necesidades iniciales}
\begin{tabularx}{\textwidth}{>{\bfseries}p{2.7cm}X X X}
\toprule
Segmento & Necesidad principal & Comportamiento o atributo & Papel en el modelo \\
\midrule
Dueño o administrador & Controlar operación, cartera, vigencias y asistencia. & Decide la compra y prefiere observar una demostración antes de definir precio. & Comprador y responsable de adopción. \\
Entrenador & Organizar planes, sesiones y seguimiento. & Utiliza información deportiva y puede asumir funciones administrativas. & Usuario profesional y posible comprador individual. \\
Miembro adulto & Consultar su plan, estado, asistencia y progreso. & Prefiere el celular y valora el acceso centralizado. & Usuario beneficiario. \\
\bottomrule
\end{tabularx}
\caption*{\fuentepropia}
\end{table}

La muestra incluyó cinco responsables y diez miembros. Los responsables aportan
evidencia sobre compra y operación; los miembros aportan evidencia sobre uso.
Esta separación se conserva en el análisis para no sumar respuestas de roles con
decisiones distintas.

La demostración debe ayudar al responsable a decidir si Pulso le facilita los
cobros y le da confianza para guardar información. Para el entrenador, debe
mostrar cómo crear y publicar una rutina sin repetir tareas. Para el miembro,
debe permitir consultar su membresía y entrenamiento desde el celular. Durante
el piloto se revisará si estas tareas se entienden y si la información se
mantiene actualizada.

\subsection{Propuesta de valor}

Pulso reúne información administrativa y deportiva en pantallas distintas según
el usuario. En el prototipo documentado, la vista del entrenador permite revisar
pagos atrasados, activar membresías y consultar asistencia. Desde Planes se
configuran las rutinas y se revisa el progreso. El miembro consulta su membresía
y sigue el entrenamiento que le asignaron.

Los principales diferenciadores son el acompañamiento inicial, el precio en
colones, las pantallas adaptables y la integración de
administración y seguimiento deportivo. Una propuesta de valor es más sólida
cuando relaciona beneficios concretos con las necesidades del segmento
\parencite{kotler2018}. La encuesta orienta Pulso hacia cobros y vencimientos
para responsables, y hacia consulta móvil de rutina y progreso para miembros.
Nueve miembros indicaron que no disponen de información digital y nueve
utilizarían principalmente el celular.

\subsection{Canales}

Los canales se organizan según la etapa del vínculo:

\begin{itemize}
  \item \textbf{Conocimiento}: demostraciones breves, redes sociales, página
  informativa y recomendaciones.
  \item \textbf{Evaluación y venta}: visita o videollamada, diagnóstico básico,
  demostración y prueba piloto.
  \item \textbf{Entrega}: sistema en línea adaptable y configuración inicial.
  \item \textbf{Soporte}: WhatsApp Business, correo y material de ayuda.
\end{itemize}

En la fase institucional se priorizó el contacto directo para conocer las
necesidades del cliente. El escenario regional incorpora un presupuesto para
publicidad pagada, que se revisará al comenzar la promoción. Dos
responsables seleccionaron soporte por WhatsApp entre los apoyos esperados y
cuatro solicitaron una demostración antes de responder sobre precio. Por ello,
demostración y contacto directo se mantienen como canales iniciales; la muestra
no permite ordenar los demás canales promocionales.

\subsection{Estrategia de mercadeo y posicionamiento}

Pulso se posiciona como un sistema en línea, accesible y acompañado para
gimnasios que necesitan ordenar su gestión sin procesos complejos. La
estrategia inicial combina los cuatro componentes de mercadeo y utiliza la
reserva operativa ya incorporada en la estructura de costos.

\begin{table}[htbp]
\caption{Estrategia inicial de mercadeo}
\begin{tabularx}{\textwidth}{>{\bfseries}p{2.8cm}X X}
\toprule
Componente & Decisión inicial & Aplicación y medición \\
\midrule
Producto & Sistema integral en línea con incorporación y soporte inicial. & Demostrar miembros, planes de entrenamiento, membresías, cobros y asistencia como recorridos diferenciados. \\
Precio & Escenario de un Starter de CRC 15\,000 y clientes Pro Gym de CRC 25\,000. & Registrar aceptación, objeciones y conversión después del piloto. \\
Distribución & Acceso web responsive y configuración directa. & Medir activaciones y uso por rol durante la prueba. \\
Promoción & Visitas, WhatsApp Business, recomendaciones y contenido demostrativo. & Contactar al menos dos gimnasios por semana durante la búsqueda del aliado y registrar cada respuesta. \\
Posición & Cobro en colones e integración administrativa y deportiva. & Comparar las razones de interés o rechazo frente a alternativas actuales. \\
\bottomrule
\end{tabularx}
\caption*{\textit{Nota}. Metas iniciales sujetas a validación; no representan resultados ya obtenidos.}
\end{table}

\subsection{Relaciones con clientes}

El equipo ayudará a configurar el sistema, resolverá consultas y revisará el
uso durante el piloto. Antes de empezar se acordarán las funciones que se
probarán, las personas responsables y la duración. Al terminar se revisará qué
se utilizó, qué problemas hubo y si el gimnasio desea continuar.

Los tutoriales y los mensajes claros ayudarán a resolver dudas frecuentes.
El equipo atenderá los problemas de cuenta, las dificultades de configuración
y las sugerencias de mejora.

Antes de cargar datos reales se revisarán los permisos y las copias de
seguridad. Durante el piloto se anotarán las consultas y el tiempo de respuesta.
La siguiente tabla resume los compromisos propuestos.

\begin{table}[htbp]
\caption{Seguimiento inicial de la calidad del servicio}
\begin{tabularx}{\textwidth}{>{\bfseries}p{3.3cm}X X}
\toprule
Etapa & Control propuesto & Evidencia esperada \\
\midrule
Incorporación & Configuración y permisos verificados antes de cargar datos reales. & Lista de comprobación aceptada por la persona responsable. \\
Operación & Registro de incidencias y respuesta inicial dentro de un día hábil durante el piloto. & Bitácora de consultas, tiempos y solución aplicada. \\
Continuidad & Revisión semanal de uso y dificultades. & Acuerdos de mejora y funciones priorizadas. \\
Cierre & Evaluación de utilidad, satisfacción e intención de pago. & Resultado de continuidad y recomendación del plan adecuado. \\
\bottomrule
\end{tabularx}
\caption*{\textit{Nota}. Los compromisos se validarán con el primer gimnasio aliado antes de convertirse en condiciones comerciales definitivas.}
\end{table}

\subsection{Fuentes de ingresos}

La fuente principal es la suscripción mensual por gimnasio. La cuota se cobra al
establecimiento y cubre el uso de responsables con rol entrenador y miembros. Como
servicios posteriores se consideran configuración especializada, soporte
prioritario, dominio propio y una aplicación independiente. Antes de contratar
un plan regular, el gimnasio puede probar la solución durante cuatro semanas por
CRC 15\,000.

\begin{table}[htbp]
\caption{Estructura comercial propuesta para Pulso}
\label{tab:planes}
\begin{tabularx}{\textwidth}{>{\bfseries}p{2.0cm}p{2.5cm}p{2.1cm}X r}
\toprule
Plan & Segmento & Límite de referencia & Alcance comercial & Precio mensual \\
\midrule
Starter & Entrenador o gimnasio pequeño. & Hasta 80 miembros. & Oferta inicial incluida en el escenario financiero; identidad visual básica prevista. & CRC 15\,000 \\
Pro Gym & Gimnasio pequeño o mediano. & Hasta 200 miembros. & Oferta principal para los clientes posteriores de la proyección; configuración visual ampliada por validar. & CRC 25\,000 \\
Premium & Gimnasio mediano o grande. & Hasta 500 miembros. & Línea posterior con creación visual propia, atención dedicada y evaluación de base aislada. & CRC 45\,000 \\
Enterprise & Cadena o caso especial. & Más de 500 miembros. & Servicio dedicado con identidad propia y posible base aislada por cadena, sujeto a cotización. & Desde CRC 70\,000 \\
\bottomrule
\end{tabularx}
\caption*{\textit{Nota}. Elaboración propia. Son suscripciones para gimnasios, con precios y capacidades por comprobar. Solo Starter y Pro Gym se incluyen en la proyección financiera. Premium y Enterprise requieren pruebas de capacidad y una cotización específica.}
\end{table}

La tarifa Starter coincide con el monto más alto del rango elegido por una
persona encuestada. Pro Gym tiene un precio propuesto para los cálculos del
negocio. Ambos se revisarán según los costos, el soporte necesario y la
aceptación de posibles clientes. Premium y Enterprise se plantean para una
etapa posterior; el punto de equilibrio de seis clientes se calcula solo con
Starter y Pro Gym.

Los planes para gimnasios grandes podrían necesitar diseño visual propio,
atención dedicada y una base de datos independiente. Su precio dependerá de la
cantidad de usuarios, el almacenamiento, la configuración y las horas de
soporte necesarias.

Cuatro responsables solicitaron personalización. Se propone comenzar con
nombre, logo y colores por gimnasio. Un dominio propio o una aplicación
independiente exigirían trabajo adicional y se estudiarían para Premium o
Enterprise. Estas funciones siguen pendientes en la oferta comercial.

\subsection{Recursos clave}

\begin{longtable}{>{\bfseries}p{3.3cm}p{11.2cm}}
\caption{Recursos clave de Pulso}\label{tab:recursos}\\
\toprule
Categoría & Recursos \\
\midrule
\endfirsthead
\toprule
Categoría & Recursos \\
\midrule
\endhead
Físicos & Computadoras disponibles, teléfonos de prueba, conexión a internet y equipo de demostración. \\
Tecnológicos & Código fuente, frontend, API, base PostgreSQL, despliegue web, repositorio, herramientas de prueba y sistema de respaldo. \\
Intelectuales & Conocimiento de programación, diseño de interfaces, bases de datos, seguridad, documentación y modelo de negocio. \\
Humanos & Dos estudiantes con responsabilidades compartidas, docente tutor, personas piloto y posibles responsables de gimnasios aliados. \\
Financieros & Aportes de ambos estudiantes, infraestructura de nube, dominio y reserva para contingencias. \\
\bottomrule
\caption*{\fuentepropia}
\end{longtable}

\subsection{Actividades clave}

Las actividades centrales son investigación, priorización del producto,
desarrollo, pruebas, despliegue, documentación, demostración, incorporación de
clientes y soporte. Las responsabilidades se comparten desde el inicio y las
decisiones técnicas y comerciales se toman de forma conjunta.

\begin{table}[htbp]
\caption{Actividades, responsables y evidencia}
\begin{tabularx}{\textwidth}{>{\bfseries}p{3.5cm}p{2.7cm}X X}
\toprule
Actividad & Responsable & Recurso & Evidencia \\
\midrule
Investigación de mercado & Equipo & Encuesta y exportación original & 15 respuestas válidas analizadas. \\
Desarrollo y corrección & Equipo & Repositorio y ambiente local & Historial y pruebas. \\
Diseño de datos y API & Equipo & Django, DRF y PostgreSQL & Modelos, endpoints y esquema. \\
Interfaz y experiencia & Equipo & React y herramientas frontend & Pantallas y pruebas. \\
Validación técnica & Equipo & Suite automatizada y demo & Reportes y compilaciones. \\
Revisión académica & Equipo y tutor & Lineamientos ExpoTÉCNICA & Lista de cumplimiento. \\
Piloto y soporte & Equipo & Gimnasio aliado & Por obtener. \\
\bottomrule
\end{tabularx}
\caption*{\fuentepropia}
\end{table}

\subsection{Alianzas clave}

El proyecto cuenta con el acompañamiento del docente tutor. Todavía se necesita
un gimnasio que acepte participar en la demostración y el piloto. Las personas
encuestadas aportaron opiniones de forma anónima, sin asumir compromisos
comerciales. Los servicios de nube se consideran proveedores.

\begin{table}[htbp]
\caption{Colaboraciones, proveedores y alianzas requeridas}
\begin{tabularx}{\textwidth}{>{\bfseries}p{3.4cm}X X p{2.3cm}}
\toprule
Relación & Aporte & Beneficio & Estado \\
\midrule
Gimnasio aliado & Contexto, prueba y retroalimentación. & Validación operacional. & Requerido; no confirmado \\
Personas encuestadas & Opiniones agregadas sobre necesidades y adopción. & Orientación inicial del producto. & Colaboración concluida; no es alianza \\
Docente tutor & Revisión académica y administrativa. & Cumplimiento. & Colaboración activa \\
Proveedores de nube & Infraestructura tecnológica contratada o gratuita. & Despliegue de la demostración. & Proveedores; no son socios \\
\bottomrule
\end{tabularx}
\caption*{\fuentepropia}
\end{table}

\subsection{Estructura de costos}

Los costos se separan en desembolsos actuales, aportes en especie e infraestructura
comercial. Esta clasificación evita afirmar que el desarrollo es gratuito cuando
utiliza equipo, internet y tiempo. El detalle y los escenarios se presentan en la
sección de costos e ingresos.
